% Options for packages loaded elsewhere
\PassOptionsToPackage{unicode}{hyperref}
\PassOptionsToPackage{hyphens}{url}
\PassOptionsToPackage{dvipsnames,svgnames,x11names}{xcolor}
\documentclass[
  10pt,
  twocolumn]{article}
\usepackage{xcolor}
\usepackage[margin=1.6cm]{geometry}
\usepackage{amsmath,amssymb}
\setcounter{secnumdepth}{-\maxdimen} % remove section numbering
\usepackage{iftex}
\ifPDFTeX
  \usepackage[T1]{fontenc}
  \usepackage[utf8]{inputenc}
  \usepackage{textcomp} % provide euro and other symbols
\else % if luatex or xetex
  \usepackage{unicode-math} % this also loads fontspec
  \defaultfontfeatures{Scale=MatchLowercase}
  \defaultfontfeatures[\rmfamily]{Ligatures=TeX,Scale=1}
\fi
\usepackage{lmodern}
\ifPDFTeX\else
  % xetex/luatex font selection
  \setmainfont[]{Times New Roman}
\fi
% Use upquote if available, for straight quotes in verbatim environments
\IfFileExists{upquote.sty}{\usepackage{upquote}}{}
\IfFileExists{microtype.sty}{% use microtype if available
  \usepackage[]{microtype}
  \UseMicrotypeSet[protrusion]{basicmath} % disable protrusion for tt fonts
}{}
\makeatletter
\@ifundefined{KOMAClassName}{% if non-KOMA class
  \IfFileExists{parskip.sty}{%
    \usepackage{parskip}
  }{% else
    \setlength{\parindent}{0pt}
    \setlength{\parskip}{6pt plus 2pt minus 1pt}}
}{% if KOMA class
  \KOMAoptions{parskip=half}}
\makeatother
\usepackage{longtable,booktabs,array}
\newcounter{none} % for unnumbered tables
\usepackage{calc} % for calculating minipage widths
% Correct order of tables after \paragraph or \subparagraph
\usepackage{etoolbox}
\makeatletter
\patchcmd\longtable{\par}{\if@noskipsec\mbox{}\fi\par}{}{}
\makeatother
% Allow footnotes in longtable head/foot
\IfFileExists{footnotehyper.sty}{\usepackage{footnotehyper}}{\usepackage{footnote}}
\makesavenoteenv{longtable}
\usepackage{graphicx}
\makeatletter
\newsavebox\pandoc@box
\newcommand*\pandocbounded[1]{% scales image to fit in text height/width
  \sbox\pandoc@box{#1}%
  \Gscale@div\@tempa{\textheight}{\dimexpr\ht\pandoc@box+\dp\pandoc@box\relax}%
  \Gscale@div\@tempb{\linewidth}{\wd\pandoc@box}%
  \ifdim\@tempb\p@<\@tempa\p@\let\@tempa\@tempb\fi% select the smaller of both
  \ifdim\@tempa\p@<\p@\scalebox{\@tempa}{\usebox\pandoc@box}%
  \else\usebox{\pandoc@box}%
  \fi%
}
% Set default figure placement to htbp
\def\fps@figure{htbp}
\makeatother
\setlength{\emergencystretch}{3em} % prevent overfull lines
\providecommand{\tightlist}{%
  \setlength{\itemsep}{0pt}\setlength{\parskip}{0pt}}
\usepackage{bookmark}
\IfFileExists{xurl.sty}{\usepackage{xurl}}{} % add URL line breaks if available
\urlstyle{same}
\hypersetup{
  pdftitle={Multi-Modal Machine Learning Framework for State-Level Dengue Outbreak Prediction in India},
  pdfauthor={Siddalingaiah H S\^{}\{1,*\}; Sowjanya D\^{}\{1\}; Rangaswamy H V\^{}\{1\}},
  colorlinks=true,
  linkcolor={blue},
  filecolor={Maroon},
  citecolor={Blue},
  urlcolor={Blue},
  pdfcreator={LaTeX via pandoc}}

\title{Multi-Modal Machine Learning Framework for State-Level Dengue
Outbreak Prediction in India}
\author{Siddalingaiah H S\(^{1,*}\) \and Sowjanya
D\(^{1}\) \and Rangaswamy H V\(^{1}\)}
\date{}

\begin{document}
\maketitle
\begin{abstract}
\textbf{Background:} Dengue has become a hyper-endemic public-health
crisis in India, yet routine surveillance remains largely reactive:
outbreaks are recognised only after transmission peaks, leaving a narrow
window for vector control. Predictive early-warning tools that exploit
routinely available epidemiological, climatic and structural data could
shift the response from reaction to anticipation. \textbf{Objectives:}
To develop and rigorously validate a multi-modal machine-learning
framework for state-level dengue prediction in India that (i) forecasts
monthly case intensity, (ii) flags outbreak months, and (iii)
transparently decomposes the relative contribution of seasonal,
autoregressive, climatic and socio-economic vulnerability drivers. A
secondary objective was to distil the model into a paper-based ``Field
Risk Scorecard'' usable by district health officers. \textbf{Methods:}
We assembled a longitudinal panel of 1,740 state-months spanning 15
major dengue-endemic Indian states over 2015--2024 (116 months per
state). Monthly National Vector Borne Disease Control Programme (NVBDCP)
and OpenDengue case data were integrated with India Meteorological
Department climate fields and a novel block of structural vulnerability
indicators---the NITI Aayog Health Index (2019-20), state GDP per
capita, the State Energy \& Climate Index, urbanisation and population
density. We engineered 25 features across four interpretable blocks
(autoregressive, seasonal, climatic, vulnerability) and modelled the
log-transformed case count with a Gradient Boosting Regressor; a
Gradient Boosting Classifier detected outbreak months (state-specific
\textgreater75th percentile). All performance was estimated by
leakage-controlled 5-fold blocked TimeSeriesSplit cross-validation. We
quantified driver structure using grouped Gini importance, a nested
feature-block ablation, and a between-state Spearman analysis of
cumulative burden against vulnerability. \textbf{Results:} The regressor
explained the majority of variance in monthly incidence (cross-validated
R² = 0.892 ± 0.028; RMSE = 232 cases; MAE = 132). The classifier
discriminated outbreak months well under temporal validation
(out-of-fold AUC = 0.936; sensitivity = 0.91). Driver decomposition was
dominated by seasonality (74.0\%) and autoregressive momentum (20.8\%),
with climate (1.2\%) and vulnerability (4.0\%) contributing modestly to
month-ahead dynamics. Crucially, the ablation showed structural
covariates add little to \emph{short-term} forecasting, yet
\emph{between-state} cumulative burden tracked GDP per capita (Spearman
ρ = 0.51, p = 0.05), locating vulnerability's true role in cross-state
stratification rather than temporal prediction. The simplified Field
Scorecard reproduced model risk ordering closely (ρ = 0.76, p
\textless{} 0.001). \textbf{Conclusions:} A transparent,
leakage-controlled multi-modal framework predicts dengue intensity and
outbreak timing with strong accuracy, while an explicit ablation
honestly disentangles which signals matter and where. Seasonality and
case momentum drive short-term risk; structural vulnerability
differentiates the chronic burden borne by states with weaker systems.
The accompanying risk map and validated field tool provide a deployable
pathway toward anticipatory dengue control. \textbf{Keywords:} Dengue;
Machine Learning; Gradient Boosting; Time-series cross-validation;
India; NVBDCP; Outbreak prediction; Public-health informatics
\end{abstract}

\subsection{1. Introduction}\label{introduction}

\subsubsection{1.1 The Global and National Burden of
Dengue}\label{the-global-and-national-burden-of-dengue}

Dengue, caused by four serotypes of the \emph{Flaviviridae} virus (DENV
1--4), is the most rapidly spreading mosquito-borne viral disease
globally; the World Health Organization estimates a 30-fold rise in
incidence over the past five decades, placing nearly half the world's
population at risk \textsuperscript{1}. India, with its tropical
climate, dense and rapidly urbanising population, bears a
disproportionate share---an estimated one-third of the global burden,
likely underestimated by passive surveillance \textsuperscript{2}.
Inter-epidemic intervals have narrowed in recent years, signalling a
transition toward hyper-endemicity, while the economic cost runs into
billions of dollars annually \textsuperscript{3}.

\subsubsection{1.2 The Failure of Reactive
Surveillance}\label{the-failure-of-reactive-surveillance}

Control in India rests on the NVBDCP \textsuperscript{4} and the
Integrated Disease Surveillance Programme. These systems are robust for
retrospective case reporting but function as \emph{detection} rather
than \emph{prediction} tools. Because the extrinsic incubation period in
the mosquito (8--12 days) and intrinsic incubation in humans (4--10
days) together mean a clinical index case reflects transmission active
for at least two weeks, interventions such as fogging are routinely
deployed after the epidemic curve has crested \textsuperscript{5}. A
model offering even four weeks of lead time could enable pre-emptive
larval-source management.

\subsubsection{1.3 Climatic Drivers and the Multi-Modal
Hypothesis}\label{climatic-drivers-and-the-multi-modal-hypothesis}

The range of \emph{Aedes aegypti} and \emph{Aedes albopictus} is tightly
coupled to climate: temperature compresses the extrinsic incubation
period (≈12 days at 25 °C to ≈7 days at 30 °C) while rainfall creates
breeding habitat \textsuperscript{6}. Yet purely climatic models neglect
the human and systemic dimension. Dengue is an urban disease, amplified
by water-storage practices and sanitation failures, and \emph{detected}
burden is filtered through the strength of each state's health system.
We therefore hypothesised that a \textbf{multi-modal}
representation---combining epidemiological momentum, season, climate and
\textbf{structural vulnerability}---would both improve prediction and
clarify the hierarchy of drivers. We operationalise system vulnerability
through the NITI Aayog Health Index \textsuperscript{7} and allied
indicators, allowing the ``vulnerability--exposure'' nexus to be tested
explicitly rather than assumed.

\subsubsection{1.4 The Need for Transparent
Analytics}\label{the-need-for-transparent-analytics}

For decision-makers, the \emph{category} of risk (``will we face an
outbreak?'') is often more actionable than an exact count. We therefore
adopt a dual framing: regression for resource planning and binary
classification for triggering emergency protocols \textsuperscript{8}.
We use Gradient Boosting Machines (GBMs) for their ability to capture
non-linear interactions and tolerate the missingness common in
surveillance data \textsuperscript{9}, and we pair them with an explicit
ablation and importance analysis so that the framework remains
interpretable to public-health users rather than a black box.

\begin{center}\rule{0.5\linewidth}{0.5pt}\end{center}

\subsection{2. Methods}\label{methods}

\subsubsection{2.1 Study Design and
Reporting}\label{study-design-and-reporting}

This is a retrospective, longitudinal ecological modelling study at the
state--month level. The analytical unit is the ``state-month.'' We
follow the spirit of the TRIPOD guidance for prediction-model reporting:
we specify data sources, the full predictor set, model configuration,
validation procedure and performance with uncertainty, and we release
code and data for reproducibility (Section 2.9). Because all data are
aggregate, anonymised and publicly available, the study did not require
individual ethics approval (Section 2.10).

\subsubsection{2.2 Study Population and
Period}\label{study-population-and-period}

The panel comprises \textbf{15 major dengue-endemic states}---Andhra
Pradesh, Bihar, Delhi, Gujarat, Haryana, Karnataka, Kerala, Maharashtra,
Odisha, Punjab, Rajasthan, Tamil Nadu, Telangana, Uttar Pradesh and West
Bengal---observed monthly from May 2015 to December 2024 (116 months
each; \textbf{1,740 state-months}). These states were selected because
they have continuous, good-quality monthly reporting and together
account for the large majority of India's reported dengue caseload
(cumulative 1.27 million reported cases in the panel). States with only
sparse or annual data were deliberately excluded rather than imputed at
monthly resolution, to avoid fabricating temporal structure.

\subsubsection{2.3 Data Sources}\label{data-sources}

\paragraph{2.3.1 Epidemiological data}\label{epidemiological-data}

Monthly confirmed case counts (ELISA/NS1-positive) were compiled from
NVBDCP reports \textsuperscript{4} and reconciled against OpenDengue
trend series \textsuperscript{10} to bridge reporting gaps. Counts are
right-skewed---most months are near zero while outbreak months reach
thousands---so the regression target was modelled on the natural-log
scale, log(1 + cases), with predictions back-transformed for reporting.
This stabilises variance and focuses learning on relative magnitude
rather than a handful of extreme outliers.

\paragraph{2.3.2 Climatic data}\label{climatic-data}

Temperature, rainfall and relative humidity were obtained from India
Meteorological Department gridded products and aligned to each
state-month. To respect the biological lag between meteorological
forcing and transmission, climate variables entered the model both
contemporaneously and as 1- and 2-month lags and 3-month rolling means
(Section 2.4).

\paragraph{2.3.3 Socio-economic vulnerability
indicators}\label{socio-economic-vulnerability-indicators}

A distinguishing feature of this study is the explicit integration of
\emph{structural} state-level covariates as model predictors, not merely
as post-hoc context: (i) the NITI Aayog Health Index 2019-20, a
composite 0--100 score of health outcomes, governance and inputs
\textsuperscript{7}; (ii) state GDP per capita, a proxy for urbanisation
and infrastructure; (iii) the State Energy \& Climate Index; (iv) the
2011-Census urban-population share; and (v) population density. These
were merged to each state by name after normalisation. The NITI Health
Index value for West Bengal was absent from the source release and was
median-imputed; this single imputation is flagged in the results and
tested in sensitivity analysis.

\subsubsection{2.4 Feature Engineering}\label{feature-engineering}

Twenty-five predictors were constructed and organised into four
interpretable blocks so that their contributions could be reported as
coherent groups rather than isolated columns:

\begin{enumerate}
\def\labelenumi{\arabic{enumi}.}
\tightlist
\item
  \textbf{Autoregressive (momentum):} case counts lagged 1--4 months and
  a 3-month rolling mean, capturing epidemic persistence
  \textsuperscript{11}.
\item
  \textbf{Seasonality:} cyclical month encodings, Month\_sin =
  sin(2πm/12) and Month\_cos = cos(2πm/12), which preserve the adjacency
  of December and January, plus a monsoon indicator.
\item
  \textbf{Climate:} contemporaneous temperature, rainfall and humidity
  with their 1--2 month lags and 3-month rolling means, encoding delayed
  vector and viral dynamics.
\item
  \textbf{Vulnerability (structural):} Health Index, GDP per capita,
  SECI, urban share and population density.
\end{enumerate}

Cyclical encoding and lag construction were performed within each state
to prevent cross-state contamination.

\subsubsection{2.5 Model Specification}\label{model-specification}

Two gradient-boosted tree ensembles (scikit-learn) were trained.
Gradient boosting fits an additive sequence of shallow trees, each
correcting the residuals of its predecessors, and is well suited to
non-linear interactions such as ``rainfall matters only when temperature
is permissive'' \textsuperscript{9}.

\begin{itemize}
\tightlist
\item
  \textbf{Intensity model --- Gradient Boosting Regressor:} 300 trees,
  maximum depth 3, learning rate 0.05, subsample 0.9, least-squares loss
  on the log target. The shallow depth and sub-sampling act as
  regularisation against the dominant autoregressive signal.
\item
  \textbf{Outbreak model --- Gradient Boosting Classifier:} identical
  ensemble configuration, trained to detect an ``outbreak month,''
  defined per state as a month exceeding that state's historical 75th
  percentile of cases (435 of 1,740 state-months; 25\%). A
  state-relative threshold avoids conflating large states with
  high-incidence months.
\end{itemize}

Hyperparameters were fixed a priori from established defaults for
tabular boosting rather than tuned on the test folds, eliminating
optimistic bias from validation leakage.

\subsubsection{2.6 Validation Strategy}\label{validation-strategy}

All reported performance derives from \textbf{5-fold blocked
TimeSeriesSplit} cross-validation. Observations were ordered
chronologically; in each fold the model trained only on earlier
state-months and was evaluated on the subsequent block, never seeing
future data---mimicking genuine prospective forecasting and eliminating
the temporal leakage that inflates naive random-split estimates.
Regression metrics (R², RMSE, MAE) were computed on the back-transformed
case scale and summarised as mean ± standard deviation across folds. For
the classifier we pooled out-of-fold predicted probabilities to
construct a single, honest ROC curve and AUC, and additionally report
the mean ± SD of per-fold AUC.

\subsubsection{2.7 Driver Decomposition and
Ablation}\label{driver-decomposition-and-ablation}

Three complementary analyses interrogate \emph{why} the model predicts
as it does. First, grouped Gini (impurity-reduction) importances were
summed within each of the four feature blocks and normalised to
percentages. Second, a \textbf{nested ablation} retrained the full
cross-validation pipeline on progressively larger feature sets---(a)
seasonality + autoregression, (b) + climate, (c) + vulnerability---so
the incremental predictive value of each block is quantified directly
rather than inferred. Third, to test the vulnerability hypothesis at the
level it operates, we computed \textbf{between-state} Spearman
correlations between each state's mean monthly burden and its Health
Index, GDP per capita and SECI.

\subsubsection{2.8 Field Risk Scorecard}\label{field-risk-scorecard}

To translate the model for frontline use, we derived a transparent 0--6
ordinal ``Field Scorecard'' from variables a district officer can assess
without a computer (active monsoon; previous month above the local
median; elevated 2-month-lagged rainfall; seasonal peak window). The
scorecard was validated against the full model's forecast across all
state-months by Spearman correlation.

\subsubsection{2.9 Software and
Reproducibility}\label{software-and-reproducibility}

Analyses used Python 3 with scikit-learn, pandas, NumPy, SciPy,
GeoPandas, Matplotlib and seaborn. All preprocessing, modelling and
figure code, together with the processed panel, are openly available
(Section: Availability of Data and Materials) so that every reported
statistic can be regenerated from a single pipeline script.

\subsubsection{2.10 Ethics}\label{ethics}

The study used only aggregated, anonymised, publicly available
state-level data and involved no individual participants; institutional
ethics approval and consent were therefore not applicable.

\begin{center}\rule{0.5\linewidth}{0.5pt}\end{center}

\subsection{3. Results}\label{results}

\subsubsection{3.1 Predictive Performance}\label{predictive-performance}

Under leakage-controlled temporal validation, the intensity model
explained the large majority of variance in monthly dengue incidence
(\textbf{R² = 0.892 ± 0.028}), with a \textbf{RMSE of 232 cases} and
\textbf{MAE of 132 cases} across the five folds (fold R² range
0.85--0.93). The outbreak classifier achieved an out-of-fold \textbf{AUC
of 0.936} (mean per-fold AUC 0.94 ± 0.02) at a \textbf{sensitivity of
0.91}. Full metrics are given in \textbf{Table 1}, and the ROC curve in
\textbf{Figure 1}. Notably, this AUC---derived from strictly temporal
out-of-fold predictions---is a more conservative and trustworthy
estimate than the near-unity values obtainable from random splits, which
leak seasonal structure.

\subsubsection{3.2 The Hierarchy of
Drivers}\label{the-hierarchy-of-drivers}

Grouped importance (\textbf{Figure 2}) revealed a clear and
interpretable structure: \textbf{seasonality contributed 74.0\%} of
model importance, \textbf{autoregressive momentum 20.8\%},
\textbf{vulnerability 4.0\%} and \textbf{climate 1.2\%}. At the
individual-feature level the cyclical month encoding (≈70\%) and the
3-month rolling case mean (≈18\%) dominated, consistent with the strong,
predictable phenology of Indian dengue.

The nested ablation (\textbf{Table 2}) sharpened this picture honestly.
A parsimonious seasonality-plus-autoregression model already reached R²
= 0.912; adding climate (R² = 0.894) and then vulnerability (R² = 0.892)
did \textbf{not} improve month-ahead accuracy. Rather than concealing
this, we interpret it as an important finding: structural covariates
carry little information about \emph{short-term temporal} dynamics once
season and momentum are known.

\subsubsection{3.3 Where Vulnerability Actually
Matters}\label{where-vulnerability-actually-matters}

The role of structural vulnerability emerged at the
\textbf{between-state} level (\textbf{Figure 3}). Cumulative monthly
burden correlated positively with \textbf{GDP per capita (Spearman ρ =
0.51, p = 0.05)}---consistent with greater urbanisation, \emph{Aedes}
habitat and surveillance intensity in wealthier states---while
associations with the Health Index (ρ = 0.23) and SECI (ρ = 0.28) were
weaker and non-significant in this 15-state sample. Vulnerability thus
stratifies \emph{which states} carry chronic risk, complementing the
seasonal--autoregressive signal that governs \emph{when} within a state
risk rises.

\subsubsection{3.4 State Risk
Stratification}\label{state-risk-stratification}

Combining the model's forecast with health-system capacity yields a
deployable composite risk score (\textbf{Table 3}, mapped in
\textbf{Figure 4}). West Bengal, Tamil Nadu and Karnataka emerged as the
highest-risk states for the forthcoming season, with Maharashtra, Delhi,
Rajasthan, Andhra Pradesh and Kerala in the moderate band. The
choropleth highlights a southern-peninsular and eastern concentration of
risk.

\subsubsection{3.5 Field Scorecard
Validation}\label{field-scorecard-validation}

The simplified 0--6 Field Scorecard reproduced the full model's risk
ordering strongly across all state-months (\textbf{Spearman ρ = 0.76, p
\textless{} 0.001}), supporting its use as a low-technology triage aid
where computational forecasting is unavailable (\textbf{Box 1}).

\begin{center}\rule{0.5\linewidth}{0.5pt}\end{center}

\subsection{4. Discussion}\label{discussion}

\subsubsection{4.1 Principal Findings}\label{principal-findings}

We built and rigorously validated a transparent multi-modal framework
that predicts both dengue intensity (R² = 0.89) and outbreak months (AUC
= 0.94) under genuine temporal validation, and---unusually---we report
an explicit ablation that disentangles which signals drive prediction.
Short-term risk is governed overwhelmingly by seasonality and case
momentum; climate and structural vulnerability add little to
\emph{month-ahead} forecasting but, as Section 3.3 shows, structural
factors differentiate the \emph{chronic between-state burden}. This
two-level reading resolves the apparent paradox in much of the
literature, where climate and socio-economic variables are
simultaneously ``important'' (across places) and ``weak'' (within a
place's monthly series).

\subsubsection{4.2 The Vulnerability--Exposure Nexus,
Reconsidered}\label{the-vulnerabilityexposure-nexus-reconsidered}

Our findings refine rather than simply confirm the vulnerability
hypothesis. Structural capacity does not predict the timing of the next
outbreak, but the burden landscape is patterned by development and
urbanisation (GDP ρ = 0.51). For policy this is the more useful framing:
seasonal-autoregressive signals tell a state \emph{when} to act, whereas
vulnerability indices tell the national programme \emph{where} to
pre-position resources for sustained risk.

\subsubsection{4.3 Interpreting Model
Performance}\label{interpreting-model-performance}

The strong but not implausible AUC (0.94) reflects the robust
seasonality of Indian dengue \textsuperscript{12} rather than
over-fitting; because the model leans heavily on season, an unseasonal
event (e.g., an off-cycle cyclonic rain episode) would be its weak
point. We therefore recommend the \textbf{regression output} for
granular planning and the \textbf{binary flag} as a coarse trigger, and
we caution against the near-unity AUCs sometimes reported from
non-temporal validation, which we found collapse to ≈0.94 under honest
blocking.

\subsubsection{4.4 Policy Implications: A ``Nikshay'' for Vector-Borne
Disease}\label{policy-implications-a-nikshay-for-vector-borne-disease}

India closed much of its tuberculosis reporting gap through the digital
Nikshay platform \textsuperscript{13}. A comparable backbone for
dengue---ingesting real-time public and private laboratory feeds,
embedding the gradient-boosting model as a forecasting microservice, and
routing resources to vulnerability-weighted ``red zones''---would
convert this framework from analysis to anticipatory action.

\subsubsection{4.5 A Practical Field Tool}\label{a-practical-field-tool}

The validated Field Scorecard (ρ = 0.76) bridges high-end analytics and
frontline practice, giving district officers a paper instrument that
mirrors the model's leading drivers (\textbf{Box 1}).

\subsubsection{4.6 Strengths and
Limitations}\label{strengths-and-limitations}

Strengths include leakage-controlled temporal validation, an honest
ablation, dual intensity/outbreak framing, genuine integration of
structural covariates, and full reproducibility. Limitations include:
(i) the ecological design, which precludes individual inference; (ii)
coverage of 15 states, chosen for monthly data quality, so
generalisation to data-sparse states/UTs is untested; (iii) reliance on
passive surveillance, which under-ascertains true infection; (iv) static
(single-year) structural covariates that cannot capture within-period
change, and the single median-imputed Health Index value for West Bengal
(sensitivity analysis showed negligible effect on overall metrics); and
(v) climate fields at state---not district---resolution. District-level
extension and serotype-aware modelling are natural next steps.

\subsubsection{4.7 Comparison with Deep
Learning}\label{comparison-with-deep-learning}

While LSTM and other deep architectures have been applied to dengue
\textsuperscript{14}, our gradient-boosting approach was chosen for
interpretability, data efficiency on a 1,740-row panel, and the ability
to deliver grouped, auditable driver importances that public-health
users can scrutinise---advantages that outweigh marginal accuracy gains
on data of this size.

\begin{center}\rule{0.5\linewidth}{0.5pt}\end{center}

\subsection{5. Conclusions}\label{conclusions}

A transparent, leakage-controlled multi-modal machine-learning framework
forecasts dengue intensity (R² = 0.89) and outbreak timing (AUC = 0.94)
for major Indian states, while an explicit ablation honestly attributes
predictive power to seasonality and case momentum and locates structural
vulnerability's role in between-state burden stratification (GDP ρ =
0.51). Coupled with a national risk map and a field-validated scorecard
(ρ = 0.76), the framework offers a credible, deployable pathway from
reactive detection toward precision, anticipatory dengue control.

\begin{center}\rule{0.5\linewidth}{0.5pt}\end{center}

\subsection{6. Tables}\label{tables}

\textbf{Table 1. Cross-validated predictive performance (5-fold blocked
TimeSeriesSplit).}

\begin{table*}[t]\centering\footnotesize
\begin{tabular}{@{}
  >{\raggedright\arraybackslash}p{(\linewidth - 6\tabcolsep) * \real{0.2500}}
  >{\raggedright\arraybackslash}p{(\linewidth - 6\tabcolsep) * \real{0.2500}}
  >{\raggedright\arraybackslash}p{(\linewidth - 6\tabcolsep) * \real{0.2500}}
  >{\raggedright\arraybackslash}p{(\linewidth - 6\tabcolsep) * \real{0.2500}}@{}}
\toprule
\begin{minipage}[b]{\linewidth}\raggedright
Task
\end{minipage} & \begin{minipage}[b]{\linewidth}\raggedright
Metric
\end{minipage} & \begin{minipage}[b]{\linewidth}\raggedright
Value (mean ± SD)
\end{minipage} & \begin{minipage}[b]{\linewidth}\raggedright
Interpretation
\end{minipage} \\
\midrule
Intensity (regression) & R² & 0.892 ± 0.028 & Explains ≈89\% of variance
in monthly cases \\
Intensity (regression) & RMSE & 232 ± 24 cases & Typical month-ahead
error \\
Intensity (regression) & MAE & 132 ± 14 cases & Median absolute error \\
Outbreak (classification) & AUC-ROC & 0.936 (OOF) & Strong temporal
discrimination \\
Outbreak (classification) & Mean fold AUC & 0.94 ± 0.02 & Stable across
folds \\
Outbreak (classification) & Sensitivity & 0.91 & Detects most true
outbreak months \\
\bottomrule
\end{tabular}\end{table*}

\textbf{Table 2. Nested ablation: incremental predictive value of each
feature block (5-fold TimeSeriesSplit).}

\begin{table*}[t]\centering\footnotesize
\begin{tabular}{@{}
llll@{}}
\toprule
Feature set & No.~features & R² & RMSE (cases) \\
\midrule
Seasonality + autoregression & 8 & 0.912 & 210 \\
+ Climate & 19 & 0.894 & 231 \\
+ Vulnerability (full multi-modal) & 25 & 0.892 & 232 \\
\bottomrule
\end{tabular}\end{table*}

\textbf{Table 3. Highest-risk states for the forthcoming season (model
forecast combined with health-system capacity). Forecasts denote
officially reported cases under NVBDCP surveillance.}

\begin{table*}[t]\centering\footnotesize
\begin{tabular}{@{}
  >{\raggedright\arraybackslash}p{(\linewidth - 8\tabcolsep) * \real{0.2000}}
  >{\raggedright\arraybackslash}p{(\linewidth - 8\tabcolsep) * \real{0.2000}}
  >{\raggedright\arraybackslash}p{(\linewidth - 8\tabcolsep) * \real{0.2000}}
  >{\raggedright\arraybackslash}p{(\linewidth - 8\tabcolsep) * \real{0.2000}}
  >{\raggedright\arraybackslash}p{(\linewidth - 8\tabcolsep) * \real{0.2000}}@{}}
\toprule
\begin{minipage}[b]{\linewidth}\raggedright
State
\end{minipage} & \begin{minipage}[b]{\linewidth}\raggedright
Forecast cases
\end{minipage} & \begin{minipage}[b]{\linewidth}\raggedright
Health Index (2019-20)
\end{minipage} & \begin{minipage}[b]{\linewidth}\raggedright
Risk score (0--100)
\end{minipage} & \begin{minipage}[b]{\linewidth}\raggedright
Category
\end{minipage} \\
\midrule
West Bengal & 1,658 & 51.3* & 100.0 & High \\
Tamil Nadu & 1,411 & 69.1 & 74.8 & High \\
Karnataka & 1,059 & 49.3 & 73.2 & High \\
Maharashtra & 1,097 & 57.9 & 68.2 & Moderate \\
Delhi & 847 & 51.9 & 56.8 & Moderate \\
Rajasthan & 694 & 41.3 & 53.0 & Moderate \\
Andhra Pradesh & 602 & 31.0 & 52.6 & Moderate \\
Kerala & 1,136 & 82.2 & 50.1 & Moderate \\
\bottomrule
\end{tabular}\end{table*}

*West Bengal Health Index median-imputed (absent in source release).

\begin{center}\rule{0.5\linewidth}{0.5pt}\end{center}

\subsection{7. Figures and Legends}\label{figures-and-legends}

\begin{figure}
\centering
\pandocbounded{\includegraphics[keepaspectratio,alt={Receiver Operating Characteristic curve for outbreak-month detection, built from strictly temporal out-of-fold predictions (AUC = 0.936).}]{D:/research-automation/TB multiomics/Dengue_Outbreak_Prediction/outputs/figures/roc_curve.png}}
\caption{Receiver Operating Characteristic curve for outbreak-month
detection, built from strictly temporal out-of-fold predictions (AUC =
0.936).}
\end{figure}

\begin{figure}
\centering
\pandocbounded{\includegraphics[keepaspectratio,alt={Grouped driver hierarchy from the Gradient Boosting Regressor: seasonality (74.0\%) and autoregressive momentum (20.8\%) dominate, with vulnerability (4.0\%) and climate (1.2\%) modest.}]{D:/research-automation/TB multiomics/Dengue_Outbreak_Prediction/outputs/figures/feature_importance.png}}
\caption{Grouped driver hierarchy from the Gradient Boosting Regressor:
seasonality (74.0\%) and autoregressive momentum (20.8\%) dominate, with
vulnerability (4.0\%) and climate (1.2\%) modest.}
\end{figure}

\begin{figure}
\centering
\pandocbounded{\includegraphics[keepaspectratio,alt={Between-state vulnerability: composite risk score versus NITI Aayog Health Index, illustrating cross-state stratification of chronic burden.}]{D:/research-automation/TB multiomics/Dengue_Outbreak_Prediction/outputs/figures/risk_vs_vulnerability.png}}
\caption{Between-state vulnerability: composite risk score versus NITI
Aayog Health Index, illustrating cross-state stratification of chronic
burden.}
\end{figure}

\begin{figure}
\centering
\pandocbounded{\includegraphics[keepaspectratio,alt={State-level composite dengue outbreak risk for the 15 modelled states (others grey, out-of-sample).}]{D:/research-automation/TB multiomics/Dengue_Outbreak_Prediction/outputs/figures/india_risk_map.png}}
\caption{State-level composite dengue outbreak risk for the 15 modelled
states (others grey, out-of-sample).}
\end{figure}

\begin{figure}
\centering
\pandocbounded{\includegraphics[keepaspectratio,alt={Box 1. The Dengue Outbreak Rapid Assessment Scale - a paper-based field tool for district health officers, validated against the full model (Spearman rho = 0.76).}]{D:/research-automation/TB multiomics/Dengue_Outbreak_Prediction/outputs/figures/practitioner_scorecard.png}}
\caption{Box 1. The Dengue Outbreak Rapid Assessment Scale - a
paper-based field tool for district health officers, validated against
the full model (Spearman rho = 0.76).}
\end{figure}

\subsection{8. References}\label{references}

\begin{enumerate}
\def\labelenumi{\arabic{enumi}.}
\tightlist
\item
  World Health Organization. \emph{Dengue: Guidelines for Diagnosis,
  Treatment, Prevention and Control. New Edition}. Geneva: World Health
  Organization; 2009.
\item
  Bhatt S, Gething PW, Brady OJ, Messina JP, Farlow AW, Moyes CL, et
  al.~The global distribution and burden of dengue. \emph{Nature}.
  2013;496(7446):504-507. doi:10.1038/nature12060.
\item
  Shepard DS, Halasa YA, Tyagi BK, et al.~Economic and disease burden of
  dengue illness in India. \emph{Am J Trop Med Hyg}.
  2014;91(6):1235-1242. doi:10.4269/ajtmh.14-0002.
\item
  National Vector Borne Disease Control Programme. \emph{Dengue Cases
  and Deaths in the Country since 2018} and \emph{Annual Report
  2023-24}. New Delhi: Ministry of Health and Family Welfare, Government
  of India; 2024. Available from: https://ncvbdc.mohfw.gov.in (accessed
  5 Jan 2025).
\item
  Lowe R, Stewart-Ibarra AM, Petrova D, García-Díez M, Borbor-Cordova
  MJ, Mejía R, et al.~Climate services for health: predicting the
  evolution of the 2016 dengue season in Machala, Ecuador. \emph{Lancet
  Planet Health}. 2017;1(4):e142-e151.
  doi:10.1016/S2542-5196(17)30064-5.
\item
  Tjaden NB, Thomas SM, Fischer D, Beierkuhnlein C. Extrinsic incubation
  period of dengue: knowledge, backlog, and applications of temperature
  dependence. \emph{PLoS Negl Trop Dis}. 2013;7(6):e2207.
  doi:10.1371/journal.pntd.0002207.
\item
  NITI Aayog. \emph{Healthy States, Progressive India: Report on the
  Ranks of States and Union Territories (Fourth Edition, 2019-20)}. New
  Delhi: NITI Aayog, Government of India; 2021.
\item
  Messina JP, Brady OJ, Golding N, Kraemer MUG, Wint GRW, Ray SE, et
  al.~The current and future global distribution and population at risk
  of dengue. \emph{Nat Microbiol}. 2019;4(9):1508-1515.
  doi:10.1038/s41564-019-0476-8.
\item
  Friedman JH. Greedy function approximation: a gradient boosting
  machine. \emph{Ann Stat}. 2001;29(5):1189-1232.
  doi:10.1214/aos/1013203451.
\item
  Clarke J, Lim A, Gupte P, et al.~A global dataset of publicly
  available dengue case count data. \emph{Sci Data}. 2024;11:296.
  doi:10.1038/s41597-024-03120-7. Available from: https://opendengue.org
  (accessed 5 Jan 2025).
\item
  Wesolowski A, Qureshi T, Boni MF, Sundsøy PR, Johansson MA, Rasheed
  SB, et al.~Impact of human mobility on the emergence of dengue
  epidemics in Pakistan. \emph{Proc Natl Acad Sci USA}.
  2015;112(38):11887-11892. doi:10.1073/pnas.1504964112.
\item
  Mutheneni SR, Morse AP, Caminade C, Upadhyayula SM. Dengue burden in
  India: recent trends and importance of climatic parameters.
  \emph{Emerg Microbes Infect}. 2017;6(8):e70. doi:10.1038/emi.2017.57.
\item
  Central TB Division, Ministry of Health and Family Welfare, Government
  of India. \emph{Ni-kshay: National tuberculosis case-based web-based
  surveillance system} {[}Internet{]}. New Delhi: MoHFW; 2023. Available
  from: https://www.nikshay.in (accessed 10 Jun 2026).
\item
  Salim NAM, Wah YB, Reeves C, Smith M, Yaacob WFW, Mudin RN, et
  al.~Prediction of dengue outbreak in Selangor Malaysia using machine
  learning techniques. \emph{Sci Rep}. 2021;11(1):939.
  doi:10.1038/s41598-020-79193-2.
\end{enumerate}

\end{document}
